\documentclass{ibetest}
\renewcommand{\arraystretch}{1.3}
\begin{document}
\testheader
\vspace{14mm}
\begin{center}
  {\sffamily\bfseries\Huge Elec-S1 \,$\cdot$\, Part 1}\\[6pt]
  {\sffamily\bfseries\LARGE Answer keys and marking schemes}\\[8pt]
  {\large Progress Tests 1.A $\to$ 3.C}\\[3pt]
  {\normalsize Semester 1 \,$\cdot$\, Academic Year 2026--2027}
\end{center}
\vspace{9mm}

{\sffamily\bfseries How to mark yourself}
\begin{itemize}[leftmargin=7mm, itemsep=3pt]
  \item Model answers are in \textcolor{keyblue}{blue}, inside the same frames as on the
        paper. Your wording may differ; what matters is the content listed in the
        \emph{marking scheme} below each exercise.
  \item Each circled item is worth the points shown. A numerical result without its unit,
        or with a wrong unit symbol, does not earn its item.
  \item Tick-box questions: 1 mark each. The ``Why'' line explains the answer and the trap
        behind the most tempting wrong option.
  \item Red boxes (\emph{Examiner's tip}) flag common traps, including the ones that cost
        marks in Test~1.
\end{itemize}

\vspace{7mm}
{\sffamily\bfseries\large Lessons from Test~1 (21 September 2026)}\par\vspace{3mm}
\begin{center}
\small
\begin{tabular}{@{}>{\raggedright\arraybackslash}p{27mm} c
                >{\raggedright\arraybackslash}p{50mm} >{\raggedright\arraybackslash}p{56mm}@{}}
\hline
\textbf{Exercise} & \textbf{Score} & \textbf{What cost marks} & \textbf{Rule from now on}\\
\hline
1.~Unit conversion & 1/1 &
  ``KV'': capital K is the kelvin. &
  Copy prefix symbols exactly: k (kilo), M (mega), m (milli), $\mu$ (micro).\\ \hline
2.~Electric current & 4/5 &
  Unit of current given as ``C\,s$^{-1}$''; charge written ``coulombs (s)''. &
  Name \emph{and} symbol: ampere (A), a base unit; coulomb (C) $=$ A\,s; second (s).\\
  \hline
3.~Circuit diagram & 0/3 &
  Counted the drawn junction dots (4 nodes); ignored the wires (4 components); 6 branches. &
  Nodes are \emph{electrically distinct} (2); each wire segment between two dots is a
  component ($4 + 6 = 10$); branches join distinct nodes (4).\\ \hline
4.~Drift velocity & 3/5 &
  Numbers put in before isolating $v$; units missing; $1/(3.2\times10^{4})$ turned into
  $3.2\times10^{-3}$. &
  $v = i/(neS)$ first, then numbers with units, powers of ten on their own:
  $v = 3.1\times10^{-5}\ \mathrm{m\,s^{-1}}$.\\ \hline
\end{tabular}
\end{center}

\vspace{6mm}
\begin{tcolorbox}[enhanced, colback=black!3, colframe=black!45, boxrule=0.4pt, arc=1mm,
  fontupper=\small]
\textbf{Counting convention used throughout.} This follows the marking of Test~1. Points joined
only by wires form a single node, and a node must join at least three conductors (a point
between two components in series is not a node). A branch runs from one node to another.
Independent node equations: $N - 1$. Independent meshes: $B - N + 1$.
\end{tcolorbox}
\end{document}
